\documentclass[border=10pt,crop=true]{standalone}
\usepackage{circuitikz}
% Shared style for 电子学一 Chapter 2 op-amp circuit figures (CircuiTikZ).
\usetikzlibrary{calc}

\ctikzset{
  bipoles/length=0.95cm,
  bipoles/thickness=1,
  resistors/scale=1,
  capacitors/scale=1,
  label distance=2.5mm,
}

\tikzset{
  opfig/.style={font=\small},
  opfigThin/.style={font=\scriptsize},
}

% Geometry (cm) — keep identical across all op-amp schematics
\def\OpInLen{1.55}    % label → input terminal
\def\OpInGap{0.08}    % terminal → wire to op amp +
\def\OpAmpWire{1.00}  % terminal side → op amp + anchor
\def\OpOutLen{1.05}   % op amp out → output terminal
\def\OpJuncL{0.50}    % v_- → feedback junction (left)
\def\OpFbDown{0.55}   % v_- straight down before R_1
\def\OpRvert{1.20}    % R_1 vertical span to ground
\def\OpInSep{0.55}    % dual-input spacing (v_1 / v_2)
\def\OpFbStub{0.40}   % follower: stub right of output before dropping
\def\OpFbClear{0.22}  % follower: below op-amp south for horizontal feedback


% 实用微分器：输入 C 串 R，反馈 R（与 §2.4.4 对偶于 §2.4.2）
\def\OpIntVmX{-1.0}
\def\OpIntVmY{-0.5}
\def\OpIntViX{-3.6}
\def\OpIntMidX{-2.3}
\def\OpIntFbY{-2.0}
\def\OpIntFbRx{1.5}
\def\OpIntVoX{2.5}
\def\OpIntPlusGnd{1.15}

\begin{document}
\begin{circuitikz}[american, scale=1.0, transform shape]
  \coordinate (OAC) at (0, 0);
  \node[op amp, noinv input down, yscale=-1](OA) at (OAC) {};

  \coordinate (VM) at (\OpIntVmX, \OpIntVmY);
  \coordinate (Mid) at (\OpIntMidX, \OpIntVmY);
  \coordinate (Vi) at (\OpIntViX, \OpIntVmY);
  \coordinate (FBL) at (\OpIntVmX, \OpIntFbY);
  \coordinate (FBR) at (\OpIntFbRx, \OpIntFbY);
  \coordinate (OutRise) at (\OpIntFbRx, 0);
  \coordinate (Vo) at (\OpIntVoX, 0);

  \draw (OA.+) -- ++(0, -\OpIntPlusGnd) node[ground, scale=0.85]{};

  \draw (OA.-) -| (VM);
  \node[circ, fill=black, inner sep=0.9pt] at (VM) {};
  \draw (VM) to[R, l=$R$, l^=] (Mid);
  \draw (Mid) to[C, l=$C$, l^=$+$] (Vi);
  \draw (Vi) to[short, o-] ++(-\OpInGap, 0) node[left, opfig]{$v_i$};

  \draw (VM) -- (FBL);
  \draw (FBL) to[R, l_=$R$, l^=] (FBR);
  \draw (FBR) -- (OutRise) -- (OA.out);

  \draw (OA.out) to[short, -o] (Vo) node[right, opfig]{$v_o$};
\end{circuitikz}
\end{document}
